\section{Introduction}
Retrieval-Augmented Generation (RAG) is a recent development combining traditional Information Retrieval and recent advances in Natural Language Processing, where a Large Language Model (LLM) is supplied with retrieved passages and asked to find the relevant information in such passages in order to answer a knowledge-grounded question \cite{NEURIPS2020_6b493230}. Since the whole point of RAG systems is to optimize the information available to the LLM by retrieving the most relevant one(s), a stress test for these systems is the one concerning multi-hop questions, where the answer depends on multiple, incremental and separate pieces of evidence \cite{trivedi2022musique}. In such context, two critical, under-studied elements are the internal ability of the LLM to assess that the supplied evidence is now \emph{sufficient}, and to judge the successful \emph{uptake} of the decisive passage into its answer computation. Prior work has shown that residual-stream activations carry signals about context-versus-parametric conflict and context reliance \cite{context-vs-memory,SpARE}, but these are typically identified on single contexts, and they typically do not analyse if and when the model's internal dynamics can predict the two above concepts in the absence of relevant parametric knowledge to answer the question.

We then use \emph{increments}, i.e. incremental additions of context, to investigate these elements. For each multi-hop question (i.e. a question requiring multiple information at different steps of the answer), we construct nested evidence contexts $C_0\subset C_1\subset\cdots$ and run a fresh, stateless forward pass on each, reading the residual stream at a fixed anchor before the answer generation. The difference between the updated and previous contexts $\Delta h_{k}=h(C_k)-h(C_{k-1})$ isolates what the added evidence context $C_k$ did to the model's state. 
We use such difference to build positive and negative examples to train linear probes for predicting \emph{sufficiency} (whether the model has gained all the relevant information to answer the question), deriving the positive examples by the \emph{decisive increments} (i.e. the question with the added relevant paragraph) and the negative ones from both the \emph{non-sufficient increments} (added information is relevant but still not enough), as well as matched \emph{distractor increments} where we substitute the correct evidence context with irrelevant ones, so as to prevent the probes overfitting on nuisances such as context length. We similarly derive training samples and fit linear probes for the other binary task of \emph{uptake} (whether the model will answer correctly, given it has gained all the relevant information), but also for \emph{correction} (answer will go from incorrect to correct, given the update), \emph{stability} (answer won't change) and \emph{revision} (answer will change) so as to separately analyse the various information recoverable from the model internal states. From the experiments, we find that sufficiency and uptake, as well as the other additional information, are all indeed recoverable from the internal states, but steering the model along the identified dimensions from each linear probe does not alter significantly models' behaviour, indicating that such dimensions are readouts (i.e. correlates) rather than causal mechanisms for the relative task. A futher baseline experiment with text-only baseline shows that sufficiency is recoverable from text alone, while uptake is not.
% At this stage, we also duplicate each true \emph{increment} (i.e. the question with the added relevant paragraph) with an increment with different types of \emph{distractors}, i.e. the question with added irrelevant or misleading paragraph; 
% given these specular increments we can then train discriminating probes which cannot overfit on nuisances such as length, but are forced to learn the directions from the internal states matching the exact conditions we are interested in (that is, when a relevant paragraph is added to the context, creating the conditions to answer the question). Furthermore, we train separate probes for identifying \emph{sufficiency} (whether the model has gained all the relevant information), \emph{uptake} (whether the model will answer correctly, given it has gained all the relevant information), but also \emph{correction} (answer will go from incorrect to correct, given the update) and \emph{stability} (answer won't change). 
% Separating these tasks allow us to achieve a clearer picture of the information that are readable from inside the model, while a separate steering experiment establishes the non-causal status of all directions identified by such probes, since steering and patching the model in such directions do not lead to significant change in the models' behavior. 
We further employ activation addition \cite{ActivationSteering} to find a true causal direction for sufficiency which can be obtained via the mean-difference between sufficient and insufficient states and, we show, acts as a gate for the model to abstain  from answering, complementing previous literature on LLMs abstention when steering \emph{confidence} directions \cite{Kumaran2026}. Finally, we show how our results hold for two additional models family and two external validation benchmarks both by re-training and transfering the learned probes. 
% Because every decisive increment (the paragraph completing the evidence) is matched with controls that add a length-matched distractor, a duplicate of already-present support, a falsified version of the decisive paragraph, or the answer string embedded in an irrelevant passage from the same penultimate state, a probe trained on these updates cannot succeed by detecting length, answer-string presence or passage relevance alone. Each increment is also labelled with the model's own behavioural transition (wrong$\to$correct, wrong$\to$wrong, correct$\to$wrong, and so on), giving separate supervision for sufficiency, uptake, correction and stability rather than one conflated ``answer-update'' target.

Our contributions are as follows.
\begin{itemize}
  \item \textbf{A matched-increment protocol} for reading evidence-sufficiency and evidence-uptake states from single forward passes before answer generation.
  \item We show that \textbf{sufficiency and uptake are linearly decodable} from residual activation of the model and that they have nearly orthogonal directions. We further show that the found directions are readouts and not causal mechanisms.
  \item We find a \textbf{causally active direction} (the mean difference between sufficient and insufficient states) which acts as a steerable abstention gate in the tested models' families.
  \item We show the \textbf{generalisability} of our method by replicating results on HotpotQA, on RAGBench's naturally retrieved documents across five domains and single-hop questions, and on a synthetic benchmark with paraphrase, partial and contradictory evidence. In this context, we build and release our own benchmark aggregating previous efforts, which we release to foster research in multi-hop reasoning and incremental evidence processing.
\end{itemize}
The picture that emerges is that evidence sufficiency and uptake are cleanly \emph{readable} before generation, that the readout of uptake is not recoverable from text, and that neither readout is a \emph{lever}: the linear probe finds a correlate of the state, while the mechanism that turns evidence into an answer is high-dimensional and content-bearing.
